% concatenated from Classification.tex
\documentclass[a4paper,11pt]{amsart}

\usepackage{amssymb,amscd,amsmath}
\usepackage[mathcal]{eucal}
\usepackage{array,float}
%\usepackage{mathtools}

%--------------------
%\usepackage{xcolor}

\usepackage{xy}
\input xy
\xyoption{all}
\usepackage{pdflscape}
\usepackage{hyperref}
\usepackage{pdflscape}
\usepackage{enumerate}
\usepackage{tikz}
\usepackage{tikz-cd}
\usepackage{relsize}%
\usetikzlibrary{decorations.pathreplacing}
\usetikzlibrary{arrows.meta, positioning}
% \usepackage[dvipsnames]{xcolor}
% \DeclareMathOperator{\lcm}{lcm}
%\usepackage[vcentermath]{youngtab}
%\usepackage[draft]{graphicx}-
\usepackage[left=3cm, right=3cm]{geometry}
\usepackage{enumerate}
\numberwithin{equation}{section}
\numberwithin{equation}{subsection}



\newtheorem{theorem}{Theorem}[section]

\newtheorem{proposition}[theorem]{Proposition}
\newtheorem{conjecture}[theorem]{Conjecture}
\newtheorem{corollary}[theorem]{Corollary}
\newtheorem{lemma}[theorem]{Lemma}
\newtheorem{claim}[theorem]{Claim}
\newtheorem{defn}[theorem]{Definition}



\newtheorem{definition}[theorem]{Definition}
\newtheorem*{remark*}{Remark}
\newtheorem{remark}[theorem]{Remark}

\newtheorem{example}[theorem]{Example}
\newtheorem{problem}[theorem]{Problem}
\newtheorem{algorithm}[theorem]{Algorithm}
\newtheorem{question}[theorem]{Question}
\newtheorem{solution}[theorem]{Solution}

%\newtheorem*{pf}{Proof}




\newtheorem{thmABC}{Fact}
\renewcommand{\thethmABC}{\Alph{thmABC}}


\newtheorem{proABC}[thmABC]{Problem}


\newcommand{\la}{{\lambda}}
\def \eb {{\mathbf e}}
\def \gb {{\mathbf g}}
\def \cb {{\mathbf c}}

\newcommand{\cH}{{\mathcal H}}
\newcommand{\cU}{{\mathcal U}}
\newcommand{\cV}{{\mathcal V}}
\newcommand{\cF}{{\mathcal F}}
\newcommand{\cQ}{{\mathcal Q}}
\newcommand{\cL}{{\mathcal L}}
\newcommand{\cG}{{\mathcal G}}
\newcommand{\cN}{{\mathcal N}}
\newcommand{\cR}{{\mathcal R}}
\newcommand{\cC}{{\mathcal C}}
\newcommand{\cP}{{\mathcal P}}
\newcommand{\cK}{{\mathcal K}}
\newcommand{\cD}{{\mathcal D}}
\newcommand{\cT}{{\mathcal T}}
\newcommand{\cM}{{\mathcal M}}
\renewcommand{\cH}{{\mathcal{H}}}
\newcommand{\ind}{{\text{Ind}}}
\newcommand{\Hom}{{\mathrm{Hom}}}
\newcommand{\rad}{\mathrm{rad}}
\newcommand{\tra}{{\mathrm{tra}}}
\newcommand{\res}{{\mathrm{res}}}
\newcommand{\PGL}{{\mathrm{PGL}}}
\newcommand{\Ho}{{\mathrm{H}}}
\newcommand{\GL}{\mathrm{GL}}
\newcommand{\vt}{{\bar{\alpha}}^{\prime}}



\newcommand{\bQ}{{\mathbb Q}}
\newcommand{\bZ}{{\mathbb Z}}
\newcommand{\bC}{{\mathbb C}}
\newcommand{\pooja}[1]{{\color{violet} #1}}
\newcommand{\sumana}[1]{{\color{red} #1}}
\newcommand{\gurleen}[1]{{\color{blue}#1}}
\newcommand{\irr}{\mathrm{Irr}}
\newcommand{\proj}{\mathrm{Proj}}


\newcommand{\mat}[9]{\begin{pmatrix}
#1 & #2 & #3 \\
#4 & #5 & #6 \\
#7 & #8 & #9
\end{pmatrix} }
\newcommand{\tr}{\mathrm{tr}}
\title[On the CLASSIFICATION OF FINITE GROUPS]{Finite Groups with nearly half as many cyclic subgroups as elements} 

\author{Vaibhav Chhajer}

\address{School of Mathematical Sciences, National Institute of Science Education and Research, An OCC of Homi Bhabha National Institute, Bhubaneswar 752050, Odisha, India}

\email{vchhajer@niser.ac.in}

\author{Sumana Hatui}
\address{School of Mathematical Sciences, National Institute of Science Education and Research, An OCC of Homi Bhabha National Institute, Bhubaneswar 752050, Odisha, India}

\email{sumanahatui@niser.ac.in}

\author{Palash Sharma}
\address{Department of Mathematical Sciences, Indian Institute of Science Education and Research Mohali, Knowledge City, Sector 81, Mohali 140 306, Punjab, India}

\email{ms22001@iisermohali.ac.in}


\begin{document}
\subjclass[2020]{20D25; 20E34}
\date{}
\keywords{finite groups, cyclic subgroups} 
\begin{abstract}
Suppose $C(G)$ denotes the set of all cyclic subgroups of a finite group $G$, and $\mathcal{O}_{2}(G)$ denotes the number of elements of order $2$ in $G$.
 In \cite{MT}, an open problem was asked to classify the groups $G$ with  $|C(G)|=|G|-r$, where $2 \leq r \leq |G|-1$. 
In this article, first  we
show that, for an odd prime $p$, there are infinitely many groups $G$ with $|C(G)|= \frac{|G|}{2}$,  $|C(G)|=\frac{|G|}{p^{q-1}}$ (for prime $q\neq p)$,  or $|C(G)|=\frac{|G|}{2}+2^{k},  k\geq 0$.
Then, we  partially answer the open question by classifying finite groups $G$ having 
$\frac{|G|}{2}-1\leq |C(G)| \leq \frac{|G|}{2}+1$ for some fix values of  $\mathcal{O}_{2}(G)$.
Finally, we provide a complete list of finite groups $G$ having $|C(G)|=\frac{|G|+(2r+1)}{2}$ for $r\geq-1$.
\end{abstract}
\maketitle


%by obtaining the best possible bound for $|C(G)|$ depending on $\mathcal{O}_{2}(G)$. 


% for the following cases. \\
% (i) $|C(G)|=\frac{|G|}{2}$ with $\mathcal{O}_{2}(G)=1,3, \frac{|G|}{2}-2, \frac{|G|}{2}-3$.\\
 % (ii) $|C(G)|=\frac{|G|}{2}-1$ with $\mathcal{O}_{2}(G)=1,\frac{|G|}{2}-3,\frac{|G|}{2}-4$.\\
 % (iii) $|C(G)|=\frac{|G|}{2}+1$  with $\mathcal{O}_{2}(G)=3,\frac{|G|}{2},\frac{|G|}{2}-1,\frac{|G|}{2}-2$.\\

 %We also prove that, if $|C(G)|=\frac{|G|}{2}+r$, then $G$ is either a $\{2,3\}$-group or a $\{2,5\}$-group.
 
%  We also prove that, if $|C(G)|=\frac{|G|}{2}+r$, \\
%  (i)   $\mathcal{O}_{2}(G)=2r-1$, $1\leq r \leq \frac{|G|}{2}$, then $G$ is a $\{2,3\}$-group. \\
%  (ii) $\mathcal{O}_{2}(G)=2r+1$, $0\leq r \leq \frac{|G|}{2}-2$, then $G$ is a $\{2,3\}$-group or $\{2,5\}$-group.\\

\section{Introduction} 
Throughout this paper, $\mathbb{Z}_{n}$, $Q_{2^{k}}$ and $M_{2^k}$ denote the cyclic group of order $n$,  the generalized quaternion group of order $2^{k}$, and  the modular maximal-cyclic group of order $2^{k}$ respectively. By SmallGroup$[i,j]$, we denote  the $j$-th group of order $i$ in GAP’s SmallGroups library. Suppose $|G:H|$ denotes the index of a subgroup $H$ in $G$, $\phi$ denotes the Euler totient function, and $\exp(G)$ denotes the exponent of a group $G$. Let $|G|$ and $|x|$ denote the order of a group $G$ and the order of an element $x \in G$, respectively.

%Let $G_{1}\rtimes G_{2}$ denotes a semidirect product of a group $G_{1}$ with another group $G_{2}$.  $\lcm(a,b)$ the least common multiple of two natural numbers $a$ and $b$.


It is always an interesting problem to classify finite groups $G$ by their certain properties. The classification of $G$ by the order of $C(G)$  received attention recently. It is easy to see that $|C(G)|= |G|$ if and only if $G$ is an elementary abelian $2$-group.
In $2015$, Tărnăuceanu \cite{MT} classified the groups with $|C(G)|=|G|-1$, imposing the following open problem.\\

\noindent \textbf{Open problem:} Describe the finite groups $G$ with $|C(G)| = |G|-r,$ where $2 \leq r \leq |G|-1$.\\

In $2016$, the answer was given for $r=2$ in \cite{MT1}. As a continuation, all finite groups with $|G|-5\leq |C(G)|\leq |G|-3$ and $2\leq|C(G)|\leq 13$ have been classified in \cite{AA,AA1,RJL,  LCQ,SR,SR1}. 
Also, the groups $G$ for certain values of the function $\alpha(G)=\frac{|C(G)|}{|G|}$ have been  studied by several authors in \cite{GS,GL,ML}, listing a few. 

% Further work on this topic can be found in \cite{R}, \cite{L}.  


Our first main result is the following. 
\begin{theorem}\label{mainresult}
Let $p$ be an odd prime. Then, there are infinitely many groups $G$ having $|C(G)|= \frac{|G|}{2}$,  $|C(G)|=\frac{|G|}{p^{q-1}},$ for prime $q\neq p$,  or $|C(G)|=\frac{|G|}{2}+2^{k},  k\geq 0$. 
 \end{theorem}



%
%In the next result, we obtain the best possible bound of  $|C(G)|$, which depends on $\mathcal{O}_{2}(G)$.
%\begin{theorem}\label{mainresult_1}
%If $\mathcal{O}_{2}(G)=|G|-r, $ $1<r<|G|$, then $$|G|-(r-2) \leq |C(G)| \leq |G|- \frac{r-1}{2}.$$ 
%
%\end{theorem}

Since there are infinitely many groups for $|C(G)|=\frac{|G|}{2}$, it is challenging  to classify the groups for this value of $|C(G)|$. Hence, in the following result, we fix the value of $\mathcal{O}_{2}(G)$ and classify the groups $G$ having $|C(G)|=\frac{|G|}{2}$. Note that, as $|G|$ is even, $\mathcal{O}_{2}(G)$ is always an odd number.

\begin{theorem}\label{mainresult_2}
Suppose $G$ is a finite group with $|C(G)|=\frac{|G|}{2}$. Then we have the followings. 
\begin{enumerate}
      \item $\mathcal{O}_{2}(G)=1$ if and only if $G\cong \mathbb Z_8, \mathbb Z_{12}, Q_{16}$ or $ \mathbb{Z}_{3}\rtimes Q_8$.
        
\item $\mathcal{O}_{2}(G)=3$ if and only if $G\cong \mathbb Z_8 \times \mathbb Z_2, M_{16}$, SmallGroup$[32,10]:$ $ Q_8 \rtimes \mathbb Z_4$, SmallGroup$[32,13]:$ $\mathbb Z_8 \rtimes \mathbb Z_4 $, SmallGroup$[32,14]:$ $\mathbb Z_8 \rtimes \mathbb Z_4$, $Q_{16}\times \mathbb{Z}_2, \mathbb Z_{12}\times \mathbb Z_2$,$ \mathbb Z_4\times (\mathbb{Z}_{3} \rtimes \mathbb{Z}_{4})$, SmallGroup$[48,12]:$ $(\mathbb{Z}_{3} \rtimes \mathbb{Z}_{4})\rtimes \mathbb{Z}_4$, SmallGroup$[48,13]:$ $\mathbb{Z}_{12}\rtimes \mathbb{Z}_4$ or $\mathbb{Z}_2\times (\mathbb{Z}_{3}\rtimes Q_{8})$.

        

        \item There is no group with $\mathcal{O}_{2}(G)=\frac{|G|}{2}-1$ or $\mathcal{O}_{2}(G)=\frac{|G|}{2}-2$.
        
        
                \item $\mathcal{O}_{2}(G)=\frac{|G|}{2}-3$ if and only if $G\cong \mathbb Z_8$.

\end{enumerate}
\end{theorem}
In the following results, we classify the groups $G$ having $|C(G)|=\frac{|G|}{2}-1$ or $ \frac{|G|}{2}+1$ for some fix values of $\mathcal{O}_{2}(G)$.
 \begin{theorem}\label{mainresult_3}
Suppose $G$ is a finite group with $|C(G)|=\frac{|G|}{2}-1$. Then we have the followings.   
\begin{enumerate} 
        \item $\mathcal{O}_{2}(G)=1$ if and only if $G\cong \mathbb Z_{10} $ or SmallGroup$[20,1]:$ $\mathbb{Z}_{5}\rtimes \mathbb{Z}_{4}$.
        

        \item There is no group with $\mathcal{O}_{2}(G)=\frac{|G|}{2}-2$ and $\mathcal{O}_{2}(G)=\frac{|G|}{2}-3$.
        
        \item $\mathcal{O}_{2}(G)=\frac{|G|}{2}-4$ if and only if $G\cong \mathbb Z_{10}$.

         
    \end{enumerate}
\end{theorem}

  
  \begin{theorem}\label{mainresult_4}
Suppose $G$ is a finite group with $|C(G)|=\frac{|G|}{2}+1$. Then we have the followings.  
\begin{enumerate} 
\item There are infinitely many groups with $\mathcal{O}_{2}(G)=1$.
\item  There is no group with $\mathcal{O}_{2}(G)=3.$
    \item $\mathcal{O}_{2}(G)=\frac{|G|}{2}$ if and only if $G\cong \mathbb{Z}_{2}$.
    \item $\mathcal{O}_{2}(G)=\frac{|G|}{2}-1$ if and only if $G\cong \mathbb{Z}_{4}$.
    \item $\mathcal{O}_{2}(G)=\frac{|G|}{2}-2$ if and only if $G\cong \mathbb{Z}_{6}$.
\end{enumerate}
\end{theorem}

\begin{theorem}\label{mainresult_5}
 Let $G$ be a group such that $|C(G)|=\frac{|G|+(2r+1)}{2}, r\geq -1$ is an integer. 
 Then \begin{enumerate}
        \item $r=-1$ if and only if $G\cong \mathbb{Z}_5$.
        
\item $r=0$ if and only if $G\cong\{e\}$ or $G$ is a group with $\exp(G)=3$.
\item For $r\geq 1$, there is no group.
    \end{enumerate}

 \end{theorem}
 
 
 The next result gives the bound on the number of prime divisors of the order of group $G$ when  $|C(G)|=\frac{|G|}{2}+m$. 
\begin{lemma}\label{bound}
Let $G$ be a group with $|C(G)|=\frac{|G|}{2}+m$, $m$ is an integer. If $p$ is a prime divisor of $|G|$, then $p\leq \mathcal{O}_2(G)+4-2m$.
\end{lemma}
It follows from the above result that, if $G$ is a group with $|C(G)|=\frac{|G|}{2}+m$, then $\mathcal{O}_2(G)>2m-3$.
Since there are infinitely many groups $G$ with $|C(G)|=\frac{|G|}{2}, \frac{|G|}{2}+2^k$ for $k \geq 0$ (by Theorem \ref{mainresult}), it is natural to ask to classify the groups $G$ by fixing $\mathcal{O}_{2}(G)$. Hence, we state the following open question.\\

\noindent \textbf{Open problem:} Classify the finite groups $G$ with $|C(G)|=\frac{|G|}{2}+ m$, for a fixed integer $m$, for the values of $\mathcal{O}_{2}(G)$, where $2m-3<\mathcal{O}_{2}(G)< |C(G)|$.

\section{prerequisites}

 First, we recall the following basic result of group theory which will be used repeatedly in this paper.
\begin{lemma}\label{normal}
    Let $G$ be a finite group and $H\leq G$, if $H$ is a unique cyclic subgroup of order $|H|$, then $H\trianglelefteq G$. 
\end{lemma}

Define a relation $\sim$ on $G$ by the following:  $a \sim b$ if $a$ and $b$ are the generators of the same cyclic subgroup of $G$. Then, $\sim$ is an equivalence relation on $G$. Let $cl(x)$ denote the equivalence class containing $x$. Now, we describe a key step, the first stair towards the proof of the main results.
\subsection{\textbf{Key step:}}\label{keystep} Let $\mathcal{O}_{2}(G)=r \neq 0$. 
Suppose $|G|=2n$ and $\{x_i\}_{i=1}^{2n}$ be the elements of $G$ such that $x_1 = e$ and $x_i^2 = e$ for $2 \leq i \leq r+1$.
Suppose $|C(G)|=k$. Then $k \geq r+1$, without loss of generality, let $x_1, x_2, \cdots,x_k$ be the generators of those $k$ cyclic subgroups of $G$.

We have $cl(x_1) = \{e\}$ and $cl(x_i) = \{x_i\}$ for $2 \leq i \leq r+1$.
Thus $|cl(x_j)| \geq 2~ $ for  $~r+2\leq j \leq k $ and so $cl(x_j)$ must contain at least one element $x_t$ for some $t\in \{k+1, \cdots, 2n\}$. 
First, we distribute exactly one element from $\{x_{k+1}, \cdots, x_{2n}\}$ to each of $cl(x_j)$, then we are left with $s=2(n-k)+r+1$ many elements which are not assigned to any class yet.

\begin{center}
\begin{tikzpicture}[every node/.style={anchor=base}]
    \node (e=x1) at (-0.3,0) {$e=x_1$};
    \node at (2.7,0.6) {\small \text{order $2$ elements}};
    \draw[decorate,decoration={brace,amplitude=5pt}] (0.8,0.3) -- (4.4,0.3);
    \node (x2) at (1,0) {$x_2$};
    \node (x3) at (2,0) {$x_3$};
    \node (dots1) at (3,0) {$\cdots$};
    \node (x1) at (4,0) {$x_{r+1}$};
    \node (xkplus2) at (5.1,0) {$x_{r+2}$};
    \node (dots2) at (6.1,0) {$\cdots$};
    \node (x2k) at (7.1,0) {$x_{k}$};
    \node (dots2) at (6.1,-0.6) {$\cdots$};
    \node (xkplus1) at (5.1,-1.2) {$x_{k+1}$};
     \node (dots2) at (6.1,-1.2) {$\cdots$};
    \node (x2kminuskplus1) at (7.2,-1.2) {$x_{2k - (r+1)}$};
    \draw[<->] (5.1,-0.2) -- (5.1,-1.0);
    \draw[<->] (7.1,-0.2) -- (7.1,-1.0);
    \node (x2kmins) at (8.8,-1.2) {$x_{2k-r}$};
    \node (dots3) at (9.8,-1.2) {$\cdots$};
    \node (x2n) at (10.8,-1.2) {$x_{2n}$};
    \node at (9.6,-1.9) {\small \text{ $s$ many elements}};
    \draw[decorate,decoration={brace,amplitude=5pt}] (11.2,-1.4) -- (8.2,-1.4);
\end{tikzpicture}
\end{center}
Now assigning the remaining $s$ elements to the equivalence classes  is equivalent to considering all the partitions of $``s"$. We can ignore those partitions which contain at least one odd number as in those cases, the no. of generators of the corresponding cyclic subgroup will be odd, a contradiction to the fact that $\phi(k)$ is even,  $\forall k\in \mathbb{N}_{\geq 3}$. 
So we consider even partitions of $s$ and assign those $s$ elements in the equivalence classes with each possibilities.
\\

\begin{theorem}\label{|G|/p}
Let $p$ be an odd prime  such that $p \nmid |G|$ or $p=2$.  Then $|C(G\times \mathbb{Z}_{p})|=2|C(G)|$.
\end{theorem}
\begin{proof}
   Recalling that the number of cyclic subgroups of a finite group $G$ can be computed as
follows  $$|C(G)|=\sum_{x\in G}\frac{1}{\phi(|x|)}.$$ Using it for $\widetilde{G}=G\times\mathbb{Z}_{p}$, we get 
\begin{equation*}
    \begin{aligned}
        |C(\widetilde{G})|&=\sum_{(x,y)\in \widetilde{G}} \frac{1}{\phi(lcm(|x|,|y|))}\\
        &= \sum_{y=0} \frac{1}{\phi(|x|)}+\sum_{\substack{y\neq 0\\ p \mathlarger{\nmid} |x|}} \frac{1}{\phi(p|x|)}+ \sum_{\substack{y\neq 0\\ p \mathlarger{\mid} |x|}} \frac{1}{\phi(|x|)}\\ &=2|C(G)|+(p-2)\sum_{\substack{x\in G \\ p \mathlarger{\mid} |x|}} \frac{1}{\phi(|x|)}.
    \end{aligned}
\end{equation*}
 When $p=2$ or $p\nmid |G|$, we have $|C(G\times \mathbb{Z}_{p})|=2|C(G)|$. 
\end{proof}

\section{Proof of main results}
In this section we prove our main results.\\


\noindent \textbf{Proof of Theorem \ref{mainresult}.}
\begin{proof}
Let $G$ be any $2$-group with $|C(G)|=\frac{|G|}{2}$.  
Then for any prime $p$, it follows from  Theorem \ref{|G|/p},
$$|C(G\times\mathbb{Z}_{p})|=2|C(G)|=|G|=\frac{|G\times\mathbb{Z}_{p}|}{p}.$$ 
 Using GAP \cite{GAP}, $M_{16}$ is a solution for $|C(G)|=\frac{|G|}{2}$. Hence, it follows that $M_{16}\times (\mathbb Z_{2})^n$ are also solutions for $|C(G)|=\frac{|G|}{2}$ for any $n\in \mathbb N$. Now,  for odd prime $p$, $M_{16}\times (\mathbb Z_{2})^n\times \mathbb Z_p$ are solutions for $|C(G)|=\frac{|G|}{p}$. 
Now, for any prime $q\neq p$, $M_{16}\times (\mathbb Z_{2})^n\times \mathbb Z_q \times \mathbb{Z}_{p^{q-1}}$ are solutions for $|C(G)|=\frac{|G|}{p^{q-1}}$ as $|C(M_{16}\times (\mathbb Z_{2})^n\times \mathbb Z_q \times \mathbb{Z}_{p^{q-1}})|=|C(M_{16}\times (\mathbb Z_{2})^n)|\cdot |C(\mathbb{Z}_q)|\cdot |C(\mathbb{Z}_{p^{q-1}})|$.
%= \frac{16\times 2^{n}}{2}\cdot 2 \cdot ((q-1)+1)=\frac{|M_{16}\times (\mathbb Z_{2})^n\times \mathbb Z_q \times \mathbb{Z}_{p^{q-1}}|}{p^{q-1}}.$ 

% first equality follows from the fact that $|C(G\times H)|=|C(G)|\cdot |C(H)|$ when $(|G|,|H|)=1$ or $H\cong \mathbb{Z}_{2}$, second equality holds from the fact $|C(\mathbb{Z}_{n})|=$ number of divisors of $n$.


Taking $G$ a group with $\exp{(G)}=3$, it is easy to check that $|C(G)|=\frac{|G|+1}{2}$. 
% by Theorem \ref{mainresult_5}. 
Consider the group $\widetilde{G}=G\times \mathbb{Z}_{2}^{k+1}$ for  $k\geq 0$. Using Theorem \ref{|G|/p} inductively, we get $$|C(\widetilde{G})|=2^{k+1}|C(G)|=2^{k}|G|+2^{k}=\frac{|\widetilde{G}|}{2}+2^{k}.$$
\end{proof}
%\noindent  \textbf{Proof of Theorem \ref{mainresult_1}.}
%\begin{proof}
%Let $\mathcal{O}_{2}(G)=|G|-r$. There are $r-1$ elements in $G$ of order greater than $2$.
%Now, we consider the even partitions of $r-1$, where each such partition represents the number of generators of the cyclic subgroups of order greater than $2$. The maximum value of $|C(G)|$ corresponds to the partition $\underbrace{{(2,2,\cdots,2)}}_{\frac{r-1}{2} \text{times}}$, and in this case $|C(G)|=|G|-\frac{r-1}{2}$. The minimum value of $|C(G)|$ corresponds to the partition $(r-1)$, in this case $|C(G)|=|G|-{(r-2)}$. Thus, we have
%$|G|-(r-2) \leq |C(G)| \leq |G|- \frac{r-1}{2}$.  
%\end{proof}




\subsection{Groups with \texorpdfstring{$|C(G)|=\frac{|G|}{2}$}{|C(G)|=|G|/2} and \texorpdfstring{$\mathcal{O}_2(G)=1$}{O2(G)=1}}\label{o2(G)=1, |G|/2}

 Let $G$ be a group of order $2n$ with $|C(G)| = \frac{|G|}{2} = n$ and $\mathcal{O}_2(G)=1$. For $n = 1, 2, 3$, it can be easily checked that $|C(G)| \neq \frac{|G|}{2}$. So,  assume that $n\geq 4$. 
 Let $\{x_i\}_{i=1}^{2n}$ be the elements of $G$ such that $x_1 = e$ and $x_2^2 = e$. 
Taking $r=1, k=n$, by our key step given in Section (\ref{keystep}), we have $s=2$. So we are left with two elements say $x_\alpha , x_\beta \in \{x_{n+1}, \cdots, x_{2n}\}$ which are not assigned to any class yet. Let, $x_\alpha, x_\beta\in cl(x_i)$ for some $i\in \{3,\cdots,n\}$.  In this case $|cl(x_i)|=4$. Hence,   $G$ must satisfy the following properties: 


 
 
$(A)$ There is a unique cyclic subgroup (say $H_0$) of order $2$
and a unique cyclic subgroup (say $N$) of order $m$ such that $\phi(m) = 4$. Thus $m\in \{5,8,10,12\}$. 
By Lemma \ref{normal}, $N\trianglelefteq G$.

$(B)$ For $1\leq i \leq {n-3}$, there are $n-3$ cyclic subgroups, say $H_i$, of order $m_i$ such that $\phi(m_i) = 2$. Thus $m_i\in \{3,4,6\}$. 

Let's call conditions $(A)$ and $(B)$ collectively as property $(P)$.
\begin{lemma}\label{lemma1}
   There is no group $G$ satisfying $(P)$ and one of the following conditions.
   \begin{enumerate}
       \item  $N \cong \mathbb{Z}_{10}$.

       \item $N \cong \mathbb{Z}_{5}$. 

        \item $N \cong \mathbb{Z}_{8}$ and $H_i \cong \mathbb{Z}_{3}$ for some $i\in \{1,\cdots,n-3\}$.

        \item $N \cong \mathbb{Z}_8$ and $H_i \cong \mathbb{Z}_{6}$ for some $i\in \{1,\cdots,n-3\}$.

         
   \end{enumerate}
\end{lemma}
\begin{proof}
\begin{enumerate}
    \item If $\mathbb{Z}_{10} \leq G$, then $\mathbb{Z}_5<\mathbb{Z}_{10}\leq G$, which contradicts $(P)$.
    
    \item If $N\cong \mathbb{Z}_{5}$, then $NH_{0}$ is a subgroup of $G$ isomorphic to $\mathbb{Z}_{10}$, a contradiction to $(P)$.

     \item  Let $H_{i}\cong \mathbb{Z}_{3}$ for some $i \in \{1,\cdots, n-3\}$ then $NH_{i}\cong \mathbb{Z}_{8}\rtimes \mathbb{Z}_{3} \cong \mathbb{Z}_{24}$, a contradiction.
    \item Suppose $H_{i}\cong \mathbb{Z}_{6}$ for some $i \in \{1,\cdots,n-3\} $. Then there is a $k \in \{1,\cdots,n-3\}$, such that $H_k < H_i$ and $H_{k}  \cong \mathbb{Z}_{3}$, which contradicts (3).
\end{enumerate}
\end{proof}

\begin{lemma}\label{Lemma2}
  If $G$ satisfies $(P)$, then the following is true.
    \begin{enumerate}[(i)]
    \item
    If $N\cong \mathbb{Z}_{8}$, then $|G|=2^{k}$ for $k\in\{3,4\}$ and $G \cong \mathbb{Z}_{8} $ or 
    $Q_{16}$.

    \item If $N\cong \mathbb{Z}_{12}$, then $|G|=2^{a}3$ for $a\in \{2,3\}$ and $G\cong \mathbb{Z}_{12}$ or $ \mathbb{Z}_{3}\rtimes Q_8$.
    \end{enumerate}

\end{lemma}
\begin{proof}
(i)    If $N\cong \mathbb{Z}_{8}$, then by Lemma \ref{lemma1} and property $(P)$, $|G|=2^{k}$ for some $k$. Assume that $k\geq 5$ then $G$ has a subgroup $H$ of order $32$. By GAP \cite{GAP},  $H$ has more than one element of order $2$, $H\cong \mathbb{Z}_{32}$ or $H\cong Q_{32}$. As $\mathbb{Z}_{32}$ and $Q_{32}$ contradict $(P)$, there is no such group $G$ of order $2^{k}$ for $k\geq 5$. So $k\in\{3,4\}$. Now using GAP \cite{GAP} we see that the only groups that satisfy $(P)$ are $\mathbb{Z}_{8}$ and $Q_{16}.$\\

(ii) First, we show that if
$H_i\cong \mathbb{Z}_3$ or $H_i\cong \mathbb{Z}_6$,
then $H_i<N$. Suppose $H_i\cong \mathbb{Z}_3$ and $H_i\cap N = \{e\}$, then $NH_i\cong \mathbb{Z}_{12}\rtimes \mathbb{Z}_3 \cong \mathbb{Z}_{36}$, a contradiction. 
    Suppose $H_i\cong \mathbb{Z}_6$, then there is a  $H_{k}\cong \mathbb{Z}_{3}$ such that $H_{k} < H_{i}$. So $H_{k} <N$.
  The result follows since $H_i\cap N$ also contains the unique element $x_2$ of order $2$. Hence, the subgroups of $G$ isomorphic to $\mathbb Z_3$ or $\mathbb Z_6$ are unique.

In this case, we have
$|G|=2^{a}3^b$. We claim that $b=1$, suppose not; say $b\geq 2$, then $G$ admits a subgroup $H$ of order $9$, then $H\cong \mathbb{Z}_{3}\times \mathbb{Z}_{3}$ or $H\cong \mathbb{Z}_{9}$. The former is not possible because there is a unique cyclic subgroup of order $3$ in $G$, and the latter is not possible because it contradicts $(P)$.
  Now, if $a\geq 5$, then $G$ has a subgroup $H$ of order $32$, which contradicts $(P)$ by the same argument used in Lemma \ref{Lemma2} $(i)$. For $1 \leq a\leq 4$, using GAP \cite{GAP}, we have either $G\cong \mathbb{Z}_{12}$ or $\mathbb{Z}_{3}\rtimes Q_8$.
  
\end{proof}





\subsection{Groups with \texorpdfstring{$|C(G)|=\frac{|G|}{2}$}{|C(G)|=|G|/2} and \texorpdfstring{$\mathcal{O}_2(G)=3$}{O2(G)=3}} \label{o2(G)=3, |G|/2} 
% Now, as $G$ is an even group so, $\mathcal{O}_{2}(G)$ cannot be even, and thus we will move to the case when $\mathcal{O}_{2}(G)=3$. 
For $n \leq 6$, it can be easily checked $|C(G)| \neq \frac{|G|}{2}$ or $\mathcal{O}_{2}(G)\neq 3$ using GAP \cite{GAP}. So, let us assume that $n\geq 7$, and $\{x_i\}_{i=1}^{2n}$ be the elements of $G$ such that $x_1 = e$ and $x_2^2=x_{3}^2=x_{4}^2 = e$. 
Taking $r=3, k=n$, by our key step given in Section \ref{keystep} we have $s=4$. So we are left with four elements, say $x_\alpha , x_\beta,x_{\gamma},x_{\delta} \in \{x_{n+1}, \cdots, x_{2n}\}$, which have not yet been assigned to any class. Thus, we have two cases.\\  

 \textbf{Case 1: $x_{\alpha},x_{\beta},x_{\gamma},x_{\delta} \in cl(x_{i})$} for some $i\in\{5,\cdots,n\}$.


 In this case, $G$ must satisfy the following $3$ conditions.

$(A_1)$ There are three cyclic subgroups of order two and a unique cyclic subgroup (say $N$) of order $m$ such that $\phi(m) = 6 \implies m\in \{7,9,14,18\}$. By Lemma \ref{normal}, $N\trianglelefteq G$.

$(B_1)$ There are $n-5$ cyclic subgroups, say $H_1,\cdots,H_{n-5}$ with orders $m_1,\cdots,m_{n-5}$ such that $\phi(m_i) = 2\implies m_i\in \{3,4,6\} ~\forall ~i\in \{1,\cdots,n-5\}$.

 Let's call conditions $(A_1)$ and $(B_1)$ collectively as property $(P_1)$.
\\

\textbf{Case 2:} $x_{\alpha},x_{\beta}\in cl(x_i)$ and $ x_{\gamma},x_{\delta} \in cl(x_{j})$ for some $i,j\in\{5,\cdots,n\}$ such that $i\neq j$.


  In this case, $G$ must satisfy the following $3$ conditions.

$(A_2)$ There are exactly three cyclic subgroups of order 2,
and exactly two cyclic subgroups (say $K_1$ and $K_2$) of order $m_1$ and $m_2$ respectively, 
such that $\phi(m_i) = 4 \implies m_i\in \{5,8,10,12\}$, for $i=1,2$.

$(B_2)$ There are $n-6$ cyclic subgroups, say $H_1,\cdots,H_{n-6}$ with orders $m_1,\cdots,m_{n-6}$ 
such that $\phi(m_i) = 2\implies m_i\in \{3,4,6\} $ $\forall i\in \{1,\cdots,n-6\}$.

 Let's call conditions $(A_2)$ and $(B_2)$ collectively as property $(P_2)$.


 \begin{lemma}\label{lemma4}
    A group $G$ satisfying $(P_1)$ cannot have a subgroup isomorphic to $\mathbb{Z}_{18},\mathbb{Z}_{14},\mathbb{Z}_{9}$ or $\mathbb{Z}_{7}$.
\end{lemma}
\begin{proof}
     If $\mathbb{Z}_{14} \leq G$ or $\mathbb{Z}_{18} \leq G$, then $\mathbb{Z}_7<\mathbb{Z}_{14}\leq G$ or $\mathbb{Z}_9<\mathbb{Z}_{18}\leq G$ respectively. Both  contradict $(P_1)$. If $N\cong  \mathbb{Z}_{9}$, then for  a cyclic subgroup of order $2$ (say $H_0$), $NH_{0}\cong \mathbb{Z}_{9}\rtimes \mathbb{Z}_{2}\cong D_{18}$ or $\mathbb{Z}_{18}$, contradict $(P_1)$. If $N \cong \mathbb{Z}_{7}$, then $NH_{0}$ is either $D_{14}$ or $\mathbb{Z}_{14}$,  contradict $(P_{1})$.
     
    \end{proof}
    
\begin{lemma}\label{lemma5}
   There is no group $G$ satisfying $(P_2)$ and one of the following conditions.
   \begin{enumerate}
\item $K_1$ or $K_2 \cong \mathbb Z_{10}$.

\item $K_1$ or $K_2 \cong \mathbb Z_{5}$.



        \item $K_1 \cong \mathbb{Z}_{8}$ and $K_2 \cong \mathbb{Z}_{12}$.


   \end{enumerate}
\end{lemma}
\begin{proof} 
$(1)$ If $K_1\cong \mathbb{Z}_{10}$, then $K_2\cong \mathbb{Z}_5$ such that $K_2 < K_1$. Taking an element, say $x_2$, of order $2$ such that $x_2\notin K_1$, we get a subgroup $K_2\langle x_2\rangle $ of order $10$, contradicting $(P_2)$.

    
    $(2)$ Suppose $K_1\cong \mathbb{Z}_5$. If $K_2\ncong \mathbb{Z}_5$, then $K_1\trianglelefteq G$. Thus, the group will contain either $D_{10}$ or $\mathbb{Z}_{10}$, a contradiction.
    

    Now assume, $K_2\cong \mathbb{Z}_5$. Then, using the previous argument, we conclude that $K_i$ are not normal in $G$, and they are conjugates of each other. 
    If any element of order $2$ lies inside $N(K_1)$, then as $K_1\trianglelefteq N(K_1)$, we will get a subgroup of order $10$,  contradicting $(P_2)$. Also note that $3 \nmid |N(K_1)|$ as we will have an element of order $15$, contradicting $(P_{2})$. Thus $N(K_1)$ must be a $5$-group. Therefore, $|N(K_1)| = 5^k$. 
    Now, if $k>1$, we will get a subgroup of order $25$ which would be either $\mathbb{Z}_5\times \mathbb{Z}_5$ or $\mathbb{Z}_{25}$, both are not possible. Hence, $|N(K_1)| = 5 \implies |G| = 2|N(K_1)| = 10$, and no such group of order $10$ satisfies $|C(G)| = \frac{|G|}{2}$.

% Suppose $K_1\cong \mathbb{Z}_5$. If $K_2\ncong \mathbb{Z}_5$ then the group will contain either $D_{10}$ or $\mathbb{Z}_{10}$ as $K_1\trianglelefteq G$ contradicting $(P_2)$. Hence $K_2\cong \mathbb{Z}_5$,  now in this case, if any one of $K_1$ or $K_2$ is normal in $G$, then previous argument will follow. So assume that both are not normal in $G$, which means that these two subgroups are conjugates of each other and moreover they completely exhaust the list of conjugates of $K_1$, thus index of normalizer of $K_1$ is then $|G:N(K_1)|=2$. If any element of order $2$ lies inside $N(K_1)$ then since $K_1\trianglelefteq N(K_1)$ we will get a subgroup of order $10$ which would be either $D_{10}$ or $Z_{10}$ contradicting $(P_2)$. Also notice that $3 \nmid |N(K_1)|$, because then we will have an element of order $15$, contradicting $(P_{2})$. Thus $N(K_1)$ must be a $5$-group. Therefore $|N(K_1)| = 5^k$. Now if $k>1$, we will get a subgroup of order $25$ which would be either $\mathbb{Z}_5\times \mathbb{Z}_5$ or $\mathbb{Z}_{25}$, the former is not possible as it contains more than two subgroups isomorphic to $\mathbb{Z}_5$ and the later is not possible as it contains an element of order $25$ contradicting $(P_2)$. Hence, $|N(K_1)| = 5 \implies |G| = 2|N(K_1)| = 10$ and no such group  $10$ which satisfies $C(G) = \frac{|G|}{2}$.
    




$(3)$ In this case, for some $i\in\{1,\cdots,n-6\}$,  we have $H_i< K_2\leq G$  such that $H_{i} \cong \mathbb{ Z}_3$. Clearly $K_1\trianglelefteq G$ and $K_1 \cap H_{i}= \{e\}$, so $K_1H_i\leq G$ with $K_1H_i\cong \mathbb{Z}_{8}\rtimes \mathbb{Z}_{3}\cong \mathbb{Z}_{24}$, a contradiction to $(P_2)$.
\end{proof} 

\begin{lemma}\label{lemma6}
   If $G$ satisfies $(P_2)$ and $K_{1}, K_{2} \cong \mathbb{Z}_{8}$, then $G$ is a $2$-group. Furthermore $G \cong \mathbb{Z}_8\times \mathbb{Z}_2, M_{16}, $ SmallGroup$[32,10]: $ $ Q_8 \rtimes \mathbb Z_4$, SmallGroup$[32,13]:$ $\mathbb Z_8 \rtimes \mathbb Z_4 $, SmallGroup$[32,14]:$ $\mathbb Z_8 \rtimes \mathbb Z_4$ or $Q_{16}\times \mathbb{Z}_2$.
\end{lemma}
\begin{proof}
  Suppose one of the $K_{i}$ is normal (say $K_1$) in $G$. If $H_i\cong \mathbb{Z}_3$ exists, then $K_1 H_i$ is a subgroup of order $24$,  a contradiction. If $K_{i}$ are not normal, then they must be conjugates of each other. In that case, $|G:N(K_1)| = 2$, and so $N(K_1)\trianglelefteq G$. If $H_i\cong \mathbb{Z}_3$ exists, then it must lie inside $N(K_1)$; otherwise, $|H_i N(K_1)|> |G|$. But $K_1H_{i}\cong \mathbb Z_{24}$,  a contradiction. Hence, $G$ must be a $2$- group. 

 For $|G| = 16$ we have two solutions $\mathbb{Z}_8\times \mathbb{Z}_2$ and $M_{16}$. For $|G| = 32$ we have four solutions which are SmallGroup$[32,10]:$ $ Q_8 \rtimes \mathbb Z_4$, SmallGroup$[32,13]:$ $\mathbb Z_8 \rtimes \mathbb Z_4 $, SmallGroup$[32,14]:$ $\mathbb Z_8 \rtimes \mathbb Z_4$, and $Q_{16}\times \mathbb{Z}_2$.  For $|G| = 64$,  there is no solution by GAP \cite{GAP}. When $|G|\geq 128$, then $G$ admits a subgroup of order $128$, by GAP \cite{GAP} any group of order $128$ contains more than three elements of order $2$, more than two $\mathbb{Z}_8$ are sitting in it or $\mathbb{Z}_{16}$ is sitting in it, giving a contradiction in each case.
\end{proof}
\begin{lemma}\label{lemma7}
Suppose $G$ satisfies $(P_2)$ and $K_{1}, K_{2} \cong \mathbb{Z}_{12}$, then $G$ is a $\{2,3\}$-group. Furthermore $G \cong \mathbb Z_{12}\times \mathbb Z_{2},$ $\mathbb Z_4\times (\mathbb{Z}_{3}\rtimes \mathbb{Z}_{4})$, SmallGroup$[48,12]:$ $(\mathbb{Z}_{3} \rtimes \mathbb{Z}_{4})\rtimes \mathbb{Z}_4$, SmallGroup$[48,13]:$ $\mathbb{Z}_{12}\rtimes \mathbb{Z}_4$ and $\mathbb{Z}_2\times (\mathbb{Z}_{3}\rtimes Q_{8})$.
\end{lemma}
\begin{proof}
   As $K_{1}, K_{2} \cong \mathbb{Z}_{12}$ so $|G|= 2^a 3^b$. Now, if $H_i\leq G$ such that $H_i\cong \mathbb Z_3$, then we claim that $H_i$ must lie inside both $K_1$ and $K_2$ and so it is unique in $G$. On contrary lets assume that $H_i$ lies outside $K_1$ then if $K_1\trianglelefteq G$ we get a subgroup $K_1H_{i}\cong \mathbb{Z}_{12}\rtimes \mathbb{Z}_{3}\cong \mathbb{Z}_{12}\times \mathbb{Z}_{3}$, which contains more than $2$ copies of $\mathbb{Z}_{12}$, a contradiction to $(P_{2})$. If $K_2$ is a conjugate of $K_1$ then $|G:N(K_1)| = 2$ and $H_i$ should lie inside $N(K_1)$ otherwise $|H_i N(K_1)|>|G|$. Since $K_1\trianglelefteq N(K_1)$, we will again get a subgroup $K_1H_{i}\cong \mathbb{Z}_{12}\times \mathbb{Z}_{3}$, contradicting $(P_{2})$. Now we claim that $b=1$, because if $b>1$, then we get a subgroup of order $9$ which would be either $\mathbb{Z}_9$ or $\mathbb{Z}_3\times \mathbb{Z}_3$, the former contradicts $(P_2)$ and the later contradicts the fact that there is a unique $\mathbb{Z}_{3}$ in $G$. Hence, $|G| = 2^a \cdot3$. 
   
   Using GAP \cite{GAP}, for $|G| = 2^3\cdot 3$, $\mathbb Z_{12}\times \mathbb Z_{2}$ is the only solution. For $|G| = 2^4\cdot 3$ we have four solutions namely $\mathbb Z_4\times (\mathbb{Z}_{3} \rtimes \mathbb{Z}_{4})$, SmallGroup$[48,12]:$ $(\mathbb{Z}_{3} \rtimes \mathbb{Z}_{4})\rtimes \mathbb{Z}_4$, SmallGroup$[48,13]:$ $\mathbb{Z}_{12}\rtimes \mathbb{Z}_4$ and $\mathbb{Z}_2\times (\mathbb{Z}_{3}\rtimes Q_{8})$. Now for $a\geq 5$, as $H_i\trianglelefteq G$, so $G$ contains a subgroup of order $2^5\cdot3$  and all groups of order $2^5\cdot 3$ have either more than three elements of order $2$ or more than two $\mathbb{Z}_{12}$ are sitting or $\mathbb{Z}_8$ is sitting inside it contradicting $(P_2)$.
\end{proof}
   
\subsection{Groups with \texorpdfstring{$|C(G)|=\frac{|G|}{2}$}{|C(G)|=|G|/2} and \texorpdfstring{$\frac{|G|}{2}-3 \leq$} \texorpdfstring{$\mathcal{O}_2(G)$} $\leq \frac{|G|}{2}-1$}\label{2.4}

\subsubsection{Groups with \texorpdfstring{$|C(G)|=\frac{|G|}{2}$}{C(G)=|G|/2} and \texorpdfstring{$\mathcal{O}_2(G)=\frac{|G|}{2}-1$}{O2(G)=|G|/2-1}}\label{2.4.1}
 In this case, we have $\frac{|G|}{2}-1$ many cyclic subgroups of order $2$, and one of identity, hence $|C(G)|= |G| \implies |G| = 0$, a contradiction.


\subsubsection{Groups with \texorpdfstring{$|C(G)|=\frac{|G|}{2}$}{|G|/2} and \texorpdfstring{$\mathcal{O}_2(G)=\frac{|G|}{2}-2$}{|G|/2-2}}\label{2.4.2}
Suppose $|G|=2n$ and $\{x_i\}_{i=1}^{2n}$ are the elements of $G$ such that $x_1 = e$ and $x_i^2= e$ $\forall i \in \{2,3,\cdots,n-1\}$. Then the cyclic subgroup generated by $x_k$ where $k\in \{n,\cdots,2n\}$ must be the same, otherwise $|C(G)|>\frac{|G|}{2}$,
contradict our hypothesis.
Thus $|x_n| \geq n+2\implies |x_n| = 2n \implies G\cong \mathbb Z_{2n}$. As there is only one element of order $2$ in $\mathbb Z_{2n}$. Hence, $n-2 = 1\implies n=3\implies G\cong \mathbb{Z}_{6}$, which is not possible.

\subsubsection{Groups with $|C(G)|=\frac{|G|}{2}$ and $\mathcal{O}_2(G)=\frac{|G|}{2}-3$} \label{2.4.3}
Let $\{x_i\}_{i=1}^{2n}$ be the elements of $G$ such that $x_1 = e$ and $x_i^2= e$ $\forall i \in \{2,3,\cdots,n-2\}$.
Suppose $x_i, 1 \leq i \leq n$ are the generators of the elements of $C(G)$.
Using the fact that,  for each divisor $k$ of $|x_{n}|$ or $|x_{n-1}|$, there are  cyclic subgroups of  order $k$ sitting inside them, we conclude that
  $|x_{n}|$ cannot be the product of two or more distinct odd primes. Therefore, 
 we are left with the following possibilities of $|x_{n-1}|$ and $|x_n|$.
\begin{center}
\begin{tabular}{ |c|c|c|c|c| } 
\hline
 & (a) & (b) & (c) &(d)\\
\hline
$|x_{n}|$ & $2p$ & $p^{2}$  & $p$ & $4$ \\ 
$|x_{n-1}|$ & $p$ & $p$ & $q,p,4$ & $4,8$\\ 
\hline
\end{tabular}
\end{center}
where $p$ and $q$ are odd primes such that $p\neq q$. 


If $\langle x_{n-1}\rangle \leq \langle x _{n} \rangle$, then $|x_n| \geq n+3\implies |x_n| = 2n \implies G\cong \mathbb Z_{2n}$. Hence, $n-3 = 1\implies n=4\implies G\cong \mathbb{Z}_{8}$. 
 Hence,  $\langle x_{n-1}\rangle \leq \langle x _{n} \rangle$ or $\langle x_{n}\rangle \leq \langle x _{n-1} \rangle$ holds if and only if $G \cong \mathbb{Z}_{8}$.
 
Case (a) and (b) are not possible because here $\langle x_{n-1}\rangle \leq \langle x _{n} \rangle$. So $G\cong \mathbb{Z}_{8}$ which leads to a contradiction as $p\nmid |G|$.

Case (c):
If $|x_{n}|=p$ and  $|x_{n-1}|=q$,  both are normal in $G$, and  $\langle x_{n} \rangle\langle x_{n-1} \rangle \cong\mathbb{Z}_{pq}$, a contradiction. 
If $|x_{n}|=p$ and  $|x_{n-1}|=p$, then $\phi(p) + \phi(p)=n+2\implies 2p-4=n \implies p\mid 4\implies p=2$, a contradiction. 
When $|x_{n}|=p$ and  $|x_{n-1}|=4$, then $\phi(p) + \phi(4)=n+2\implies p-1=n \implies p\nmid n$, a contradiction. 

Case (d): If $|x_{n}|=|x_{n-1}|=4$, then $\phi{(4)}+\phi{(4)}=n+2\implies n=2$. Thus $G\cong \mathbb{Z}_4$, which is not possible. Finally, if $|x_{n}|=4$ and $|x_{n-1}|=8$, then $\langle x_{n}\rangle \leq \langle x_{n-1}\rangle$ thus $G\cong \mathbb{Z}_{8}$, which is not possible.\\

 




\noindent \textbf{Proof of Theorem \ref{mainresult_2}.}
\begin{proof}
 (1)  Proof follows from Section \ref{o2(G)=1, |G|/2}.\\
 (2) Proof follows from Section \ref{o2(G)=3, |G|/2}.\\
Proof of (3) and (4) follows from Section \ref{2.4}.
\end{proof}


\subsection{Groups with \texorpdfstring{$|C(G)|=\frac{|G|}{2}-1$}{|C(G)|=|G|/2-1} and \texorpdfstring{$\mathcal{O}_2(G)=1$}{O2(G)=1}}\label{o2G=1, |G|/2-1} 
 For $n \leq 4$ it can be easily checked that $|C(G)| \neq \frac{|G|}{2}-1$ using GAP \cite{GAP}. So, let us assume that $n\geq 5$. Let $\{x_i\}_{i=1}^{2n}$ be the elements of $G$ such that $x_1 = e$ and $x_2^2 = e$. 
 Taking $r=1, k=n-1$,  by our key step given in Section (\ref{keystep}) we have $s=4$. So we are left with four elements say $x_\alpha , x_\beta,x_{\gamma},x_{\delta} \in \{x_{n}, \cdots, x_{2n}\}$ which are not assigned to any class yet. Thus, we have two cases.\\  
  \textbf{Case 1: $x_{\alpha},x_{\beta},x_{\gamma},x_{\delta} \in cl(x_{i})$} for some $i\in\{3,\cdots,n-1\}$.


In this case, $G$ must satisfy the following conditions.

$(A_3)$ There is a unique cyclic subgroup of order two, and there is a unique cyclic subgroup (say $N$) of order $m$ such that $\phi(m) = 6 \implies m\in \{7,9,14,18\}$. By Lemma \ref{normal}, $N\trianglelefteq G$.

$(B_3)$ There are $n-4$ cyclic subgroups say $H_1,\cdots,H_{n-4}$ with orders $m_1,\cdots,m_{n-4}$ such that $\phi(m_i) = 2\implies m_i\in \{3,4,6\} \forall i\in \{1,\cdots,n-4\}$.


In this case, using the proof of Lemma \ref{lemma4}, we can show that there are no groups.\\ 

 \textbf{Case 2:} $x_{\alpha},x_{\beta}\in cl(x_i)$ and $ x_{\gamma},x_{\delta} \in cl(x_{j})$ for some $i,j\in\{3,\cdots,n-1\}$ such that $i\neq j$.


  In this case, $G$ must satisfy the following  conditions.

$(A_4)$ There is a unique cyclic subgroup of order two.

$(B_4)$ There are exactly two cyclic subgroups (say $K_1$ and $K_2$) of order $m_1$ and $m_2$
such that $\phi(m_i) = 4 \implies m_i\in \{5,8,10,12\}$, for $i=1,2$.

$(C_4)$ There are $n-5$ cyclic subgroups say $H_1,\cdots,H_{n-5}$ with orders $m_1,\cdots,m_{n-5}$ 
such that $\phi(m_i) = 2\implies m_i\in \{3,4,6\} $ $\forall i\in \{1,\cdots,n-5\}$.

 Let's call conditions $(A_4)$, $(B_4)$ and $(C_4)$ collectively as property $(P_4)$.

 \begin{lemma}\label{lemma8}
     There is no group $G$ satisfying $(P_4)$ and one of the following conditions:
   \begin{enumerate}
\item $K_1\cong \mathbb{Z}_{10}$ and $K_2 \ncong \mathbb Z_{5}$.
%(10,8),(10,10),(10,12)
\item  $K_1 \cong \mathbb Z_{5}$ and  $K_2 \ncong \mathbb Z_{10}$.
%(5,5),(5,8),(5,12)


        \item $K_1,K_2\in \{\mathbb{Z}_{8},\mathbb{Z}_{12}\}$


   \end{enumerate}
 \end{lemma}
\begin{proof}
   (1) and (2) follow using the similar idea given in the proof of Lemma \ref{lemma5}.

    (3) When $K_1\cong \mathbb{Z}_{8}$ and $K_2\cong \mathbb{Z}_{12}$, then result follows from the proof of Lemma \ref{lemma5}.
    
    When $K_1,K_2\cong \mathbb{Z}_8$, following the same idea from the proof of Lemma \ref{lemma6} we get that $|G| = 2^{a}$ where $a\geq 4$.  Also, $a\leq 4$ follows from the ideas in proof of Lemma \ref{Lemma2}. Thus, $a=4$. Using GAP \cite{GAP}, no group of order $16$ satisfies the hypothesis.  
    
    When $K_1,K_2\cong \mathbb{Z}_{12}$, again following the same idea from the proof of Lemma \ref{lemma7} we get that $|G| = 2^{a}3$ where $a\in \{3,4\}$. By GAP \cite{GAP}, there is no group of orders $24$ and $48$ satisfying the hypothesis.
\end{proof}
 
\begin{lemma}\label{lemma9}
    If $G$ satisfies $(P_4)$, then $G$ is a $\{2,5\}$-group. Furthermore $G \cong \mathbb Z_{10}$ or SmallGroup$[20,1]:$ $\mathbb{Z}_{5}\rtimes \mathbb{Z}_{4}$.
\end{lemma}
\begin{proof} 
By Lemma \ref{lemma8} we have $K_1\cong \mathbb{Z}_{10}$ and $K_2 \cong \mathbb Z_{5}$.
    It is also easy to check that $H_{i}\ncong \mathbb{Z}_{3},\mathbb{Z}_6$ $\forall i \in \{1,\cdots, n-5\}$, and hence $|G|=2^{a}5^{b}$ for some $a,b$. We claim that $b=1$, as if $b>1$ then we get a subgroup of order $25$ which would be either $\mathbb{Z}_{25}$ or $\mathbb{Z}_5\times \mathbb{Z}_5$ both contradicting $(P_4)$. Hence, $|G| = 2^a \cdot 5$. Also as $|G|\geq 10$, so $a\geq 1$. Following the reasoning of Lemma \ref{Lemma2}, we have $a\leq 4$. Hence $|G|\in \{10,20,40,80\}$ and using GAP \cite{GAP} we get the result.
\end{proof}

\subsection{Groups with \texorpdfstring{$|C(G)|=\frac{|G|}{2}-1$}{|C(G)|=|G|/2} and \texorpdfstring{$\frac{|G|}{2}-4 \leq$} \texorpdfstring{$\mathcal{O}_2(G)$} $\leq \frac{|G|}{2}-2$}\label{2.6}

\subsubsection{Groups with \texorpdfstring{$|C(G)|=\frac{|G|}{2}-1$}{C(G)=|G|/2-1} and \texorpdfstring{$\mathcal{O}_2(G)=\frac{|G|}{2}-2$}{O2(G)=|G|/2-3}} There is no  such group, which can be proved using similar ideas as in  Section \ref{2.4.1}.



\subsubsection{Groups with \texorpdfstring{$|C(G)|=\frac{|G|}{2}-1$}{C(G)=|G|/2-1} and \texorpdfstring{$\mathcal{O}_2(G)=\frac{|G|}{2}-3$}{O2(G)=|G|/2-3}} Following similar ideas as in Section \ref{2.4.2}, we get $G\cong \mathbb{Z}_{8}$, which is not a solution.

\subsubsection{Groups with \texorpdfstring{$|C(G)|=\frac{|G|}{2}-1$}{C(G)=|G|/2-1} and \texorpdfstring{$\mathcal{O}_2(G)=\frac{|G|}{2}-4$}{O2(G)=|G|/2-4}} 

Let $\{x_i\}_{i=1}^{2n}$ be the elements of $G$ such that $x_1 = e$ and $x_i^2= e$ $\forall i \in \{2,3,\cdots,n-3\}$.
Suppose $x_i, 1 \leq i \leq n-1$ are the generators of the elements of $C(G)$.
Using the fact that  for each divisor $k$ of $|x_{n-1}|$ or $|x_{n-2}|$ there are  cyclic subgroups of  order $k$ sitting in them, we conclude that
  $|x_{n-1}|$ cannot be product of two or more distinct odd primes, and so
 we are left with the following possibilities of $|x_{n-2}|$ and $|x_{n-1}|$.
\begin{center}
\begin{tabular}{ |c|c|c|c|c| } 
\hline
 & (a) & (b) & (c) &(d)\\
\hline
$|x_{n-2}|$ & $2p$ & $p^{2}$  & $p$ & $4$ \\ 
$|x_{n-1}|$ & $p$ & $p$ & $q,p,4$ & $4,8$\\ 
\hline
\end{tabular}
\end{center}
where $p$ and $q$ are odd primes such that $p\neq q$. 


If $\langle x_{n-1}\rangle \leq \langle x _{n-2} \rangle$, then $|x_{n-2}| \geq n+4\implies |x_{n-2}| = 2n \implies G\cong \mathbb Z_{2n}$. Hence, $n-4 = 1\implies n=5\implies G\cong \mathbb{Z}_{10}$. 
 Hence,  $\langle x_{n-1}\rangle \leq \langle x _{n-2} \rangle$ or  $\langle x_{n-2}\rangle \leq \langle x _{n-1} \rangle$ holds if and only if $G \cong \mathbb{Z}_{10}$.
 
Case (a): $\phi(2p)+\phi(p)=n+3\implies 2p-5=n\implies p=n=5\implies G\cong \mathbb{Z}_{10}$ which is a solution.

Case (b):  Here  $\langle x_{n-1}\rangle \leq \langle x _{n-2} \rangle$ holds thus $G\cong \mathbb{Z}_{10}$, but it doesn't have cyclic subgroup with order $p^{2}$ for any $p$.

Case (c):
If $|x_{n-2}|=p$ and  $|x_{n-1}|=q$ or $4$, then similar ideas as in Section \ref{2.4.3} will give a contradiction. If $|x_{n-1}|=|x_{n-2}|=p$, then $\phi(p) + \phi(p)=n+3\implies 2p-5=n\implies p=n=5$ and $|G|=10$ but both groups of order $10$ do not have two cyclic subgroups of order $5$. 

Case (d): When $|x_{n-1}|=|x_{n-2}|=4$, we have $\phi{(4)}+\phi{(4)}=n+3\implies n=1$. Thus $G\cong \mathbb{Z}_2$, which is not a solution. And finally when $|x_{n-2}|=4$ and $|x_{n-1}|=8$, then $\langle x_{n-2}\rangle \leq \langle x_{n-1}\rangle$ thus $G\cong \mathbb{Z}_{10}$, which doesn't have cyclic subgroup of order $4$ and $8$. 
\\


\noindent\textbf{Proof of Theorem \ref{mainresult_3}}
\begin{proof}
 (1) Proof follows from Section \ref{o2G=1, |G|/2-1}. 
 
 Proof of (2) and (3) follows from  Section \ref{2.6}.
\end{proof}

%\subsection{Groups with \texorpdfstring{$|C(G)|=\frac{|G|}{2}+1$}{|C(G)|=|G|/2} and \texorpdfstring{$\mathcal{O}_2(G)=n-3$}{O2(G)=n-3} or \texorpdfstring{$n-4$}{n-4}}


\noindent\textbf{Proof of Theorem \ref{mainresult_4}}
\begin{proof}
 (1) Follows from Theorem  \ref{mainresult} when $k=0$.

(2), (3), (4) and (5) can be proved using similar ideas as used in proof of Theorem \ref{mainresult_2} and Theorem \ref{mainresult_3}.
\end{proof}



\noindent \textbf{Proof of Theorem \ref{mainresult_5}}.
\begin{proof}
 In this case, the group must be of odd order. Let $\{x_i\}_{i=1}^{2n+1}$ be the elements of $G$ such that $x_1 =e$.
 
 (1)  It can be easily checked that there are no solutions for $n=0,1$. So assume that $n>1$. Then  $cl(x_k)\geq 2 $ $\forall k\in \{2,\cdots, n\}$. Therefore, repeating the same process as in the proof of Theorem \ref{mainresult_2}, we get the following conditions on $G$.


$(A_5)$ There is a unique cyclic subgroup (say $N$) of order $m$ such that $\phi(m) = 4 \implies m\in \{5,8,10,12\}$. By Lemma \ref{normal}, $N\trianglelefteq G$.

$(B_5)$ There are $n-2$ cyclic subgroups say $H_1,\cdots,H_{n-2}$ with orders $m_1,\cdots,m_{n-2}$ such that $\phi(m_i) = 2\implies m_i\in \{3,4,6\} \forall i\in \{1,\cdots,n-2\}$.

Let's call conditions $(A_5)$ and $(B_5)$ collectively as property $(P_5)$.

Since, the group is of odd order, therefore in $(A_5)$, $m=8, 10, 12$ are not possible. For the same reason, in $(B_5)$, $m_i = 4,6$ are not possible. So the only possibility is that $N\cong \mathbb Z_5$ and $H_i\cong \mathbb Z_3$ $\forall i$. Now if $n>2$ then such a $H_i$ exists and we get $NH_i\leq G$ where $NH_i\cong \mathbb Z_{15}$ contradicting $(P_5)$. So, the only possibility is $n= 2$. For $n=2$, $G\cong \mathbb Z_5$ and it satisfies $|C(G)| = \frac{|G|-1}{2}$.\\


(2) Let $|G| = 2n+1$ where $n\geq 0$. Then $G$ have $n$ cyclic subgroups say $H_1,\cdots,H_{n}$ of orders $m_1,\cdots,m_{n}$ such that $\phi(m_i) = 2\implies m_i\in \{3,4,6\} \forall i\in \{1,\cdots,n\}$.  Since $G$ is of odd order, $m_i\neq 4,6$. So every non-trivial element of $G$ must be of order $3$. Hence, $G\cong\{e\}$ or $G$ is a group with $\exp(G)=3$.\\

(3) For $r\geq 1$, repeating the same process as in the proof of (1), this time we are left with $2r$ non-trivial elements, each of them having order $m$ such that $\phi(m)=1$, hence each of them is an involution which contradicts the fact that group is odd.
\end{proof}


\noindent \textbf{Proof of Lemma \ref{bound}}.
\begin{proof}
In this case, by our key step given in \ref{keystep}, we have $s=\mathcal{O}_{2}(G)+1-2m$. Hence, the largest number of generators  a cyclic subgroup of $G$ can have $\mathcal{O}_2(G)+3-2m$ elements. For any prime divisor $p$ of  $|G|$, we have $\phi(p)=p-1 \leq \mathcal{O}_2(G)+3-2m\implies p\leq \mathcal{O}_2(G)+4-2m$. 
\end{proof}

% \section{Some more results}
Using the similar ideas used in the proof of the main results, we have the following. 
\begin{lemma}
     Let $G$ be a group of order $2n$ and $|C(G)|=n+r$ for a fixed positive integer $r$.
     \begin{enumerate}
         \item If  $1\leq  r\leq n $ and $\mathcal{O}_{2}(G)=2r-1$,  then $G$ is a $\{2\}$-group or $\{2,3\}$-group. 
         \item If $1\leq r\leq n -3 $ and $\mathcal{O}_{2}(G)=2r+1$, then $G$ is a $\{2\}$-group or $\{2,3\}$-group or $\{2,5\}$-group.
     \end{enumerate}
      \end{lemma}
%      \begin{proof}
%         (1) As $r\leq n$ so let $x_{1},x_{2},\cdots, x_{2r},\cdots,x_{n+r}$ be generators of those $n+r$ cyclic subgroups of $G$ such that $x_{1}=e$ and $x_{2}^2=\cdots=x_{2r}^{2}=e$. Repeating the key step given in Section \ref{keystep}, we get that $G$ must satisfy the followings.

%  $(A_6)$ There are $2r-1$ cyclic subgroups of order 2.

%   $(B_6)$ There are $n-r$ cyclic subgroups say $H_1,\cdots,H_{n-r}$ with orders $m_1,\cdots,m_{n-r}$ such that $\phi(m_i) = 2\implies m_i\in \{3,4,6\} \forall i\in \{1,\cdots,n-r\}$.

% Hence $G$ is a $\{2,3\}$-group or $\{2\}$-group.\\

% (2)  As $r\leq n-3$ so let $x_{1},x_{2},\cdots, x_{2r+2},\cdots,x_{n+r}$ be generators of those $n+r$ cyclic subgroups of $G$; such that $x_{1}=e$ and $x_{2}^2=\cdots=x_{2r+2}^{2}=e$. So, $G$ satisfy the followings.

% $(A_7)$ There are $2r+1$ cyclic subgroup of order 2 and there is a unique cyclic subgroup (say $N$) of order $m$ such that $\phi(m) = 4 \implies m\in \{5,8,10,12\}$. By Lemma \ref{normal}, $N\trianglelefteq G$.

% (B_7)$ There are $n-r-3$ cyclic subgroups say $H_1,\cdots,H_{n-r-3}$ with orders $m_1,\cdots,m_{n-r-3}$ such that $\phi(m_i) = 2\implies m_i\in \{3,4,6\} \forall i\in \{1,\cdots,n-r-3\}$.

% Observe that $N\cong \mathbb{Z}_{10}$ is not possible. If $N\cong \mathbb{Z}_{5}$, then $m_{i}=4$, $\forall i\in \{1,\cdots,n-r-3\}$. Hence $G$ is either $\{2,3\}$-group or $\{2,5\}$-group or $\{2\}$-group.
%      \end{proof}
%
%\noindent By Theorem \ref{mainresult_1}, it follows that to classify groups $G$ by the value of $\mathcal{O}_{2}(G)$ is same as to classify groups $G$ by some values of $|C(G)|$.
%In the following result, we give a list of groups $G$ having $\mathcal{O}_{2}(G)=|G|-r$ for $r\in\{1,3,5,7\}.$
%
%\begin{theorem}\label{ord 2 ele in G}
%    Let $G$ be a group with $\mathcal{O}_{2}(G)=|G|-r$ then 
%    \begin{enumerate}
%        \item $r=1$ if and only if $G\cong \mathbb{Z}_{2}^{n}$, for $n\geq 0$.
%        \item $r=3$ if and only if $G\cong \mathbb{Z}_{3},\mathbb{Z}_{4}, S_3$ or $D_8 $.
%        \item $r=5$ if and only if $G\cong \mathbb{Z}_{6},\mathbb{Z}_2\times\mathbb{Z}_{4}, D_{12}, \mathbb{Z}_{2}\times D_{8}, \mathbb{Z}_{5}$ or $D_{10} $.
%        \item $r=7$ if and only if $G \cong Q_8,\mathbb{Z}_{8}, D_{16}, \mathbb{Z}_{7}$ or $D_{14}$.
%    \end{enumerate}
%\end{theorem}
%
%
%\begin{proof}
%    (1) By  \cite[Theorem 1]{MT}, $G\cong \mathbb{Z}_{2}^{n}$ for $n\geq 0$.  
%
%(2) By Theorem \ref{mainresult_1}, we have $|C(G)|=|G|-1$, so by  \cite[Theorem 2]{MT} possibilities of $G$ are $\mathbb{Z}_{3},\mathbb{Z}_{4}, S_3,D_8$ and all of them satisfy $\mathcal{O}_{2}(G)=|G|-3$.
%
%(3) By Theorem \ref{mainresult_1},  we have $|C(G)|=|G|-2$ or $|G|-3$. Now, the result follows by checking the groups given in  \cite[Theorem 1]{MT1} and  \cite[Theorem 3]{RJL}.
%
%%possibilities of $G$ are $\mathbb{Z}_{6},\mathbb{Z}_2\times\mathbb{Z}_{4}, D_{12}, \mathbb{Z}_{2}\times D_{8},Q_8, \mathbb{Z}_{5}, D_{10}$.  Then it is easy to check that  $G \cong$ $\mathbb{Z}_{6},\mathbb{Z}_2\times\mathbb{Z}_{4}, D_{12}, \mathbb{Z}_{2}\times D_{8}, \mathbb{Z}_{5}, D_{10}$. 
%
%(4) By Theorem \ref{mainresult_1}, we have  $|C(G)|=|G|-3$, $|G|-4$ or $|G|-5$. Then, the result follows by checking the groups given in \cite[Theorem 3, Theorem 4, Theorem 5]{RJL}.
%
%% possibilities of $G$ are $Q_8, \mathbb{Z}_{5}, D_{10}, \mathbb{Z}_{4}\times \mathbb{Z}_{2}\times \mathbb{Z}_{2}, \mathbb{Z}_{2}\times \mathbb{Z}_{2}\times D_8, (\mathbb{Z}_{2}\times \mathbb{Z}_{2})\rtimes \mathbb{Z}_{4}, Q_{8}\rtimes\mathbb{Z}_{2}, \mathbb{Z}_{3}\times\mathbb{Z}_{3}, (\mathbb{Z}_{3}\times\mathbb{Z}_{3})\rtimes \mathbb{Z}_{2}, A_{4}, \mathbb{Z}_{6}\times\mathbb{Z}_{2}, \mathbb{Z}_{2}\times\mathbb{Z}_{2}\times S_3, \mathbb{Z}_{8}, D_{16}, \mathbb{Z}_{7}, D_{14}, \mathbb{Z}_{3}\rtimes \mathbb{Z}_{4}$, and out of them those groups which indeed satisfy $\mathcal{O}_{2}(G)=|G|-7$ are $Q_8,\mathbb{Z}_{8}, D_{16}, \mathbb{Z}_{7}, D_{14}$.
%\end{proof}
%Finally, we ask a natural open question related to the above study.
%\vspace{0.3cm}
%
%\noindent \textbf{Open problem:} Classify the finite groups $G$ with $\mathcal{O}_{2}(G)=|G|-r$, for  $  9 \leq r \leq  |G|-1$.
%

\vspace{1cm}


\noindent\textbf{Acknowledgements:}
% The authors are thankful to the editor for providing valuable suggestions which has greatly improved the presentation of the paper. We are also grateful to the anonymous referee for carefully and critically reading the manuscript and providing many valuable comments. 
The first named author acknowledges the research support of the SERB-SRG, Department of Science and Technology (Grant no. SRG/2023/000894), Govt. of India. The third named author thank the National Institute of Science Education and Research, Bhubaneswar, for providing a summer intern fellowship. The authors thank the National Institute of Science Education and Research, Bhubaneswar, for providing an excellent research facility.

\begin{thebibliography}{999}
\bibitem{AA}
A.R. Ashrafi, and E. Haghi, On the number of cyclic subgroups in a finite group. Southeast Asian Bull. Math. 42 (2018), 865–873.

\bibitem{AA1}
A.R. Ashrafi, and E. Haghi, On n-cyclic groups. Bulletin of the Malaysian Mathematical Sciences Society 42 (2019): 3233-3246.

 \bibitem{RJL}
R. Belshoff, J. Dillstrom, and  L. Reid, Finite groups with a prescribed number of cyclic subgroups. Communications in Algebra 47.3 (2019): 1043-1056.



\bibitem{LCQ}
X. Chen, H. Liu, and S. Qiao, Finite groups with a small number of cyclic subgroups. Czechoslovak Mathematical Journal (2025): 1-13.

\bibitem{GS}
X. Gao, and R. Shen, Finite groups with many cyclic subgroups. Indian Journal of Pure and Applied Mathematics 54.2 (2023): 485-498.

\bibitem{GL}
M. Garonzi, and I. Lima, On the number of cyclic subgroups of a finite group. Bulletin of the Brazilian Mathematical Society, New Series 49 (2018): 515-530.

\bibitem{ML}  
M.S. Lazorec, and M. Tărnăuceanu, A note on the number of cyclic subgroups of a finite group." Bulletin mathématique de la Société des Sciences Mathématiques de Roumanie 62.4 (2019): 403-416.



% \bibitem{R}
% I. M. Richards, A remark on the number of cyclic subgroups of a finite group. The American Mathematical Monthly. 91 (1984), 571-572.


\bibitem{SR}
K. Sharma, and A.S. Reddy, Groups having $11 $ or $12 $ cyclic subgroups. arXiv preprint arXiv:2210.11788 (2022).

\bibitem{SR1}
K. Sharma, and A.S. Reddy, Groups having $11 $ cyclic subgroups. International Journal of Group Theory 13.2 (2024): 203-214.


\bibitem{MT}
 M. Tărnăuceanu, Finite groups with a certain number of cyclic subgroups. The American Mathematical Monthly 122.3 (2015): 275-276.

\bibitem{MT1}
M. Tărnăuceanu, Finite Groups with a Certain Number of Cyclic Subgroups II. Acta Universitatis Sapientiae Mathematica 10.2 (2018): 375-377.


% \bibitem{L}
% L. Tóth: On the number of cyclic subgroups of a finite Abelian group. Bull. Math. Soc. Sci. Math. Roum., Nouv. Sér. 55 (2012), 423-428.

 \bibitem{GAP}
 The GAP Group, Groups algorithms and programming. Version 4.14.0, (2024), http://www.gap-system.org.
\end{thebibliography}








\end{document}

